%!TeX root = main.tex
%!encoding = UTF-8 Unicode

\section{Oversampling and Averaging}
The previous section explains the challenges to measure high-frequency components of a signal.
Firstly, classical digital noise reduction can be applied, which address the thermal noise \mi{v}{noise,therm}: Oversampling and averaging.

By oversampling and averaging afterwards, the effective number of bits (ENOB) can be increased by \cite{silicon_laboratories_an118_2013}.

\begin{equation}
	\mti{N}{ENOB} = \mti{N}{ADC} + 0.5\cdot \log_2(\frac{\mi{f}{s}}{2\mi{f}{max}})
\end{equation}

Here \mi{N}{ADC} is the nominal ADC resolution, \mi{f}{s} is the sampling frequency of the oscilloscope and \mi{f}{max} the maximum frequency of interest.
Furthermore, the factor of 2 in the denumerator is due to the nyquist sampling theorem.
This can be done internally in the oscilloscope.

As this cannot be pushed limitless due to hardware-restrictions of the maximum sampling speed of modern oscilloscopes, the number of signal points can be increased by not measuring one single pulse, but a pulse chain of $M$ pulses.
This further increases the ENOB to \cite{kragl_active_2025}

\begin{equation}
	\mti{N}{ENOB} = \mti{N}{ADC} + 0.5\cdot \log_2(\frac{\mi{f}{s}}{2\mi{f}{max}}) + 0.5\log_2(M)
\end{equation}

These gains are realized only if the dominant noise is random, uncorrelated with the signal, and greater than the quantization step \cite{silicon_laboraties_an118_2013}. 
As discussed above, in power electronics the high‑frequency voltage components often fall below or near the ADC quantization level; therefore oversampling and averaging alone are insufficient and must be combined with a high‑pass measurement method introduced next.
